\documentclass[a4paper,11pt]{article}

%page geometry
\usepackage[colorlinks,linkcolor=blue,citecolor=blue,urlcolor=blue]{hyperref}
\usepackage[lmargin=2cm,rmargin=2cm,tmargin=1.25cm,bmargin=1.25cm,includefoot,includehead]{geometry}

%typography
\usepackage[T1]{fontenc}
\usepackage[utf8]{inputenc}
%% Serif Times fonts
\renewcommand{\rmdefault}{ptm} 
%% Sans-serif Arial-like fonts (Helvetica)
\renewcommand{\sfdefault}{phv} 

%text structures
\usepackage{enumitem}
\usepackage{tabu}
\usepackage[dvipsnames]{xcolor}
\usepackage{tcolorbox}
\usepackage{colortbl}
\usepackage{graphicx}
\usepackage{eurosym}
\usepackage{xspace}

%uncomment if you want to draw a Gantt diagram
%\usepackage{pgfgantt}



%couleurs ARN 
\definecolor{ANRblue}{HTML}{004d95}
\definecolor{ANRred}{HTML}{FF0000}
\definecolor{ANRpurple}{HTML}{800080}


%physique-chimie
\usepackage{siunitx} %unités en \SI{10}{\micro\metre}
\usepackage[version=3]{mhchem} % Formula subscripts using \ce{}

\usepackage{csquotes}

%commandes pour dire qui doit faire quoi
\newcommand{\Mathieu}[1]{\textcolor{orange}{#1}}
\newcommand{\Helene}[1]{\textcolor{purple}{#1}}

% %biblio
\usepackage[backend=biber,style=numeric-comp, language=british,eprint=false, url=false, doi=false, sortcites=true, sorting=none, isbn=false, giveninits=true,maxbibnames=99]{biblatex}
\renewbibmacro{in:}{%
  \ifentrytype{article}{}{%
  \printtext{\bibstring{in}\intitlepunct}}}
% \bibliography{Mendeley}

%no month
\AtEveryBibitem{\clearfield{month}}

% \addbibresource{IEEEabrv.bib}

%lien DOI ou URL si disponible
\newbibmacro{string+doi}[1]{%
  \iffieldundef{doi}{%
  	\iffieldundef{url}{#1}{\href{\thefield{url}}{#1}}}{\href{http://dx.doi.org/\thefield{doi}}{#1}}}

%sur le titre, en couleur
\DeclareFieldFormat
  [article,inbook,incollection,inproceedings,patent,thesis,unpublished]
  {title}{\usebibmacro{string+doi}{\textcolor{ANRblue}{#1\addperiod}}}
  
%Our names in bold
\usepackage{xstring}
\usepackage{etoolbox}
\newboolean{bold}
%% New biblatex 3.3 way
% \newcommand{\makeauthorsbold}[1]{%
%   \DeclareNameFormat{author}{%
% 	\setboolean{bold}{false}%
% 	\renewcommand{\do}[1]{\ifstrequal{##1}{####1}{\setboolean{bold}{true}}{}}%
% 	\docsvlist{#1}%
%   \nameparts{##1}%
% 	\ifnumequal{\value{listcount}}{1}%
%     {\ifnumequal{\value{liststop}}{1}%Single author
%       {\expandafter\ifthenelse{\boolean{bold}}{\textcolor{ANRpurple}{\underline{\namepartfamily\addcomma\addspace\namepartgiveni,\isdot}}}{\namepartfamily\addcomma\addspace\namepartgiveni\addcomma\isdot}}
%       %first author
%       {\expandafter\ifthenelse{\boolean{bold}}{\textcolor{ANRpurple}{\underline{\namepartfamily\addcomma\addspace\namepartgiveni}}}{\namepartfamily\addcomma\addspace\namepartgiveni}}}
%       %last author
%     {\ifnumless{\value{listcount}}{\value{liststop}}
%       {\expandafter\ifthenelse{\boolean{bold}}{\addcomma\addspace\textcolor{ANRpurple}{\underline{\namepartfamily\addcomma\addspace\namepartgiveni}}}{\addcomma\addspace\namepartfamily\addcomma\addspace\namepartgiveni}}
%       %middle author
%           {\expandafter\ifthenelse{\boolean{bold}}{\addcomma\addspace \textcolor{ANRpurple}{\underline{\namepartfamily\addcomma\addspace \namepartgiveni\addcomma\isdot}}}{\addcomma\addspace\namepartfamily\addcomma\addspace\namepartgiveni\addcomma\isdot}}%
%     }%
%   }%
% }

%% Legacy biblatex way
% \newcommand{\makeauthorsbold}[1]{%
% \DeclareNameFormat{author}{%
% 	\setboolean{bold}{false}%
% 	\renewcommand{\do}[1]{\ifstrequal{##1}{####1}{\setboolean{bold}{true}}{}}%
% 	\docsvlist{#1}%
% 	\ifnumequal{\value{listcount}}{1}%
%     {\ifnumequal{\value{liststop}}{1}%Single author
%       {\expandafter\ifthenelse{\boolean{bold}}{\textcolor{ANRpurple}{\underline{##1\addcomma\addspace ##4\addcomma\isdot}}}{##1\addcomma\addspace ##4\addcomma\isdot}}
      %first author
%       {\expandafter\ifthenelse{\boolean{bold}}{\textcolor{ANRpurple}{\underline{##1\addcomma\addspace ##4}}}{##1\addcomma\addspace ##4}}}
      %last author
%     {\ifnumless{\value{listcount}}{\value{liststop}}
%       {\expandafter\ifthenelse{\boolean{bold}}{\addcomma\addspace \textcolor{ANRpurple}{\underline{##1\addcomma\addspace ##4}}}{\addcomma\addspace ##1\addcomma\addspace ##4}}
      %middle author
%           {\expandafter\ifthenelse{\boolean{bold}}{\addcomma\addspace \textcolor{ANRpurple}{\underline{##1\addcomma\addspace ##4\addcomma\isdot}}}{\addcomma\addspace ##1\addcomma\addspace ##4\addcomma\isdot}}%
%     }%
% }%
% }
% \makeauthorsbold{Cottin-Bizonne,Delanoë-Ayari,Leocmach}

%%hack fullcite so that it prints all authors
\makeatletter
\DeclareCiteCommand{\fullcite}
  {\defcounter{maxnames}{\blx@maxbibnames}%
    \usebibmacro{prenote}}
  {\usedriver
     {\DeclareNameAlias{sortname}{default}}
     {\thefield{entrytype}}}
  {\multicitedelim}
  {\usebibmacro{postnote}}
\makeatother


\newcolumntype{P}[1]{>{\raggedright}p{#1}}

%% section title format, in the text and in the table of content
\usepackage{tocloft}

\renewcommand{\cfttoctitlefont}{\color{ANRblue}\normalfont\scshape\bfseries\sffamily\Large}

\usepackage{titlesec}
\titleformat{\section}
{\color{ANRpurple}\normalfont\scshape\bfseries\sffamily\Large}
{\color{ANRpurple}\thesection}{1em}{}
\renewcommand{\cftsecfont}{\color{ANRpurple}\normalfont\scshape\bfseries\sffamily\large}

\titleformat{\subsection}
{\color{ANRblue}\normalfont\scshape\bfseries\sffamily\large}
{\color{ANRblue}\thesubsection}{1em}{}
\renewcommand{\cftsubsecfont}
{\color{ANRblue}\normalfont\scshape\bfseries\sffamily}

\titleformat{\subsubsection}
{\color{ANRpurple}\normalfont\bfseries\sffamily}
{\color{ANRpurple}\thesubsubsection}{1em}{}
\renewcommand{\cftsubsubsecfont}
{\color{ANRpurple}\normalfont\bfseries\sffamily\small}

\titleformat{\paragraph}[runin]
{\color{ANRblue}\normalfont\sffamily}
{\color{ANRblue}\theparagraph}{1em}{}

%titre, acronyme, défi, instrument
\newcommand{\mytitle}[0]{Very cool project}
\newcommand{\myacro}[0]{VeCoolPro}
\newcommand{\mydefi}[0]{3}
\newcommand{\myinstrument}[0]{JCJC}
\author{Coordinateur: Mathieu Leocmach\\
\small Univ Lyon, Université Claude Bernard Lyon 1, CNRS, Institut Lumière Matière
}

%% Headers
\usepackage{fancyhdr}
\pagestyle{fancy}
%\fancyhfoffset[re]{4cm}
\lhead{\myacro\\ Défi \mydefi{} - \myinstrument}
\rhead{\LARGE\textcolor{ANRblue}{AAPG ANR 2017}}
\renewcommand{\headrulewidth}{0pt}

\lfoot{}
\cfoot{}
\rfoot{\thepage}%/\pageref{LastPage}} 




\title{\vspace{-\baselineskip}\mytitle}
\date{}

\newcommand{\content}[1]{\emph{#1}\\} 


\begin{document}
\maketitle
\thispagestyle{fancy}
%\section{Contexte, positionnement et objectif(s) de la proposition}
\section{Context, position and objectives of the proposal}
%Sur une page et demie (1,5) environ, 
%Décrire les objectifs et les hypothèses scientifiques
%Montrer l’originalité et la pertinence par rapport à l’état de l’art
%Décrire la méthodologie et la gestion des risques

%\section{Organisation du projet et moyens mis en œuvre}
\section{Project organisation and implementation}

%Sur une demi (0,5) page environ,
%S’il s’agit d’un projet collaboratif (PRC, PRCE), 
%Présenter succinctement le coordinateur scientifique et son implication :
%Indiquer l’expérience du coordinateur scientifique et son rôle dans le projet
%Présenter le consortium et sa complémentarité : 
%Indiquer les différentes compétences envisagées pour mener le projet objet de la pré-proposition, en précisant l’identité du ou des scientifiques impliqués, l’identification des établissements auxquels ils sont rattachés et 2-3 références pertinentes dans le domaine en lien direct avec la pré-proposition (publications, faits marquants de R&D, brevets, prix scientifiques, produits, procédés, licences, services…), et tout autre élément permettant de juger de la qualité des partenaires et du consortium le cas échéant. 
%Décrire succinctement les moyens mis en œuvre pour atteindre les objectifs :
%Indiquer notamment le montant total de l’aide prévisionnelle demandée, sa répartition entre partenaires et entre grands postes de dépenses (équipement, personnel, prestations, autre fonctionnement) et les moyens de recherche complémentaires au projet.
%
%S’il s’agit d’un projet jeune chercheuse – jeune chercheur (JCJC), 
%Présenter succinctement le coordinateur scientifique, sa situation au moment de la soumission et lors de la mise en œuvre du projet, son rôle dans le projet 
%Présenter l’équipe du coordinateur scientifique et son positionnement scientifique au sein de l’organisme ou du laboratoire d’accueil.
%Décrire succinctement les moyens mis en œuvre pour atteindre les objectifs :
%Indiquer notamment le montant total de l’aide prévisionnelle demandée, sa répartition entre grands postes de dépenses (équipement, personnel, prestations, autre fonctionnement) et les moyens de recherche complémentaires au projet.

%\section{Impact et retombées du projet}
\section{Project impact and benefits}
%Sur une (1) page environ,
%Décrire dans quel(s) domaine(s) (scientifique, économique, social ou culturel) les résultats du projet peuvent avoir un impact
%Décrire en quoi les attendus du projet répondent aux enjeux d’un (axe de) défi du Plan d’Action 2017 de l’ANR, de la SNR, de la recherche fondamentale
%Décrire la stratégie de diffusion et/ou de valorisation envisagée : l’objectif est d’indiquer comment les résultats seront diffusés et/ou valorisés (publications scientifiques ou grand public, brevets, etc.), et non de détailler les modalités de partage de l’éventuelle propriété intellectuelle entre partenaires

%%%% Couper le PDF ici %%%%
\pagebreak
\section*{\bibname}
\printbibliography[heading=none]

\end{document}